\documentclass[uplatex,dvipdfmx,a4paper,11pt]{jsarticle}
\input{../../../templates/kaisetsu.tex}
\kaisetsuprint{熱物理学　第10章「完全な熱力学関数と Legendre 変換」}{2026-11-20}
%% kaisetsu-moto	problems\熱物理学\shou10\熱物理学-10.tex	kiroku\kaisetsu\熱物理学\熱物理学-10_2026-11-20_komento.tex
\begin{document}
\matome
\begin{saitenichiran}
  \hyougyou{\hyperlink{bt-m-077-u-tv}{2}}{L1・説明}{○◎○}{B}{(1)記述 (3)記述}
  \hyougyou{\hyperlink{bt-m-078-legendre}{3}}{L1・計算}{◎◎◎△}{A}{(4)計算}
  \hyougyou{\hyperlink{bt-m-079-power-law}{4}}{L2・導出}{◎◎○◎}{A}{(3)記述}
  \hyougyou{\hyperlink{bt-m-080-gibbs-duhem}{5}}{L2・計算}{◎○◎}{A}{(2)記述}
\end{saitenichiran}
\begin{daiji}
  \item $f$ を $x$ 方向に $C$ ずらすと、Legendre 変換は $g_{C}(y)=g(y)-Cy$ と1次の項が加わる。上下の平行移動ではない。傾き $y$ の接線を右に $C$ ずらすと切片が $-Cy$ 変わる、と図で覚える（問題3 (4)）。
  \item $p=T(\partial p/\partial T)_{V,n}$ から $p=Tf(V,n)$ へは、$(\partial/\partial T)(p/T)_{V,n}=0$ を経由する。偏微分方程式を解いて任意関数が残ることが、$U(T,V)$ が完全でないことの中身である（問題2 (1)）。
  \item 計算した条件・量が何から来るかを一言添える：$\mu\dd n$ は物質の授受に伴う仕事（問題2 (3)）、$C_{V}>0$・$\kappa_{S}>0$ は $U$ の凸性（問題4 (3)）、$\rho$ は示強変数で $n$ によらない（問題5 (2)）。
\end{daiji}
\begin{maemo}
  \item 式変形の根拠を書かない：第3章 問題5（2026-11-12）で係数比較の根拠（$\dd T$ と $\dd p$ が独立）が無かった。今回も問題2 (1) で $p=T(\partial p/\partial T)$ から $p=Tf(V,n)$ へ1行で飛んでいる。「$\therefore$」の前後で何を使ったかを1語でも書く。
  \item 小問の後半の書き落とし：第3章 問題2 (3)（2026-11-12）、第8章 問題5 (3)（2026-11-16）の $C_{V}$ に続き、今回は問題2 (3) の「$\mu$ を $\delta W$ の中の係数として説明せよ」が無い。問題文の「求めよ」「説明せよ」の数と答えの数を合わせる。
\end{maemo}

\youtenseiri
\begin{noteshikilist}
  \item その偏微分 % 出典: p.49
    \[ \left(\pdv{U}{S}\right)_{V} = T(S,V), \qquad \left(\pdv{U}{V}\right)_{S} = -p(S,V) \notetag{10.2} \]
  \item $U(T,V)$ の体積微分（エネルギー方程式） % 出典: p.49
    \[ \left(\pdv{U}{V}\right)_{T} = T\left(\pdv{p}{T}\right)_{V,n} - p \notetag{10.3} \]
  \item $U(S,V,n)$ の偏微分 % 出典: p.50
    \[ \left(\pdv{U}{S}\right)_{V,n} = T, \qquad \left(\pdv{U}{V}\right)_{S,n} = -p, \qquad \left(\pdv{U}{n}\right)_{S,V} = \mu \notetag{10.4} \]
  \item 物質の出入りがある場合の仕事 % 出典: p.50
    \[ \delta W = -p\dd V + \mu\dd n \notetag{10.5} \]
  \item 定積比熱の正値性 % 出典: p.53
    \[ C_{V} = T\left(\pdv{S}{T}\right)_{V,n} = T\Big/\left(\pdv{T}{S}\right)_{V,n} > 0 \notetag{10.9} \]
  \item 断熱圧縮率の正値性 % 出典: p.53
    \[ \kappa_{S} = -\frac{1}{V}\left(\pdv{V}{p}\right)_{S,n} = -1\Big/V\left(\pdv{p}{V}\right)_{S,n} > 0 \notetag{10.10} \]
  \item Euler の関係式 % 出典: p.53
    \[ U = TS - pV + \mu n \notetag{10.12} \]
  \item Gibbs-Duhem の式 % 出典: p.54
    \[ S\dd T - V\dd p + n\dd\mu = 0 \notetag{10.13} \]
  \item 等温圧縮率を密度で表した式 % 出典: p.54
    \[ \kappa_{T} = -\frac{1}{V}\left(\pdv{V}{p}\right)_{T,n} = \frac{1}{\rho}\left(\pdv{\rho}{p}\right)_{T} = \frac{1}{\rho^{2}}\left(\pdv{\rho}{\mu}\right)_{T} \notetag{10.15} \]
  \item Legendre 変換 % 出典: p.55
    \[ g(y) = f(x(y)) - yx(y) \notetag{10.16} \]
  \item Legendre 変換の導関数 % 出典: p.55
    \[ x(y) = -\dv{g}{y} = -g'(y) \notetag{10.17} \]
  \item 逆 Legendre 変換（$x = -\dd g/\dd y$ を $y$ について解いたものを $y = y(x)$ とする） % 出典: p.56
    \[ f(x) = g(y(x)) + xy(x) \notetag{10.18} \]
\end{noteshikilist}

\mondaihen
\level{1}
\mondaibangou{2}
\begin{mondai}{077-u-tv}{説明}{$U(T,V)$ が完全でないこと（自然な変数の組），化学的仕事と共役な変数の組，第8章のエネルギー方程式を使う}
\begin{teisuu}
  気体定数 & $R$ & \qty{8.314}{\joule\per\kelvin\per\mole}\\
\end{teisuu}
物質量 $n$ の一成分流体の内部エネルギーが $U(T,V,n) = ncT$（$c$ は正の定数）と分かっているとする。
  \begin{shoumon}
    \item エネルギー方程式 \noteshiki{10.3} を使い，この $U(T,V,n)$ と矛盾しない状態方程式は $p = Tf(V,n)$（$f$ は任意の関数）の形であることを示せ。また，$p = nRT/V$ と，$b$ を正の定数とした $p = nRT/V + bn^{2}T/V^{2}$ がどちらも矛盾しないことを確かめよ。
    \item (1) から，$U(T,V)$ だけでは状態方程式が決まらないことが分かる。一方，$U(S,V)$ が分かれば状態方程式が決まる理由を述べよ。
    \item 物質の出入りがある場合の準静的過程の仕事 $\delta W$ と，$\dd U$ に現れる共役な変数の組を3つ挙げよ。また，化学ポテンシャルを $\delta W$ の中の係数として説明せよ。
  \end{shoumon}

\begin{kotaeran}
  \machigai{1}{○}{記述} $0=T(\partial p/\partial T)_{V,n}-p$ から、すぐに「$\therefore p=Tf(V,n)$」とした。間をつなぐ議論が無い。$V$、$n$ を固定して $(\partial/\partial T)(p/T)=(1/T)(\partial p/\partial T)_{V,n}-p/T^{2}=0$ なので、$p/T$ は $T$ によらない $V$、$n$ だけの任意の関数 $f(V,n)$ になる、と書く。2つの状態方程式の確かめは正しい。
  \machigai{3}{○}{記述} $\delta W=-p\dd V+\mu\dd n$ と共役な組 $(T,S)$、$(p,V)$、$(\mu,n)$ は正しい。問われている「化学ポテンシャルを $\delta W$ の中の係数として説明せよ」の答えが無い。$\mu\dd n$ は物質の授受に伴う仕事（化学的仕事）で、$\mu$ は $S$ と $V$ を一定に保って物質量を $\dd n$ だけ増やすときに系に与えられる仕事の係数である。
\end{kotaeran}
\begin{rikai}
  \item $T(\partial p/\partial T)=p$ は「$V$、$n$ を固定すると $p\propto T$」という意味である。積分定数にあたるものが $V$、$n$ の任意関数として残るので、$U(T,V)$ だけでは状態方程式が1つに決まらない。
  \item $-p\dd V$ が体積を変える仕事であるのと同じく、$\mu\dd n$ は物質を出し入れする仕事である。$\dd U=T\dd S+\delta W$ の $\delta W$ の各項が「示強変数 $\times$ 示量変数の変化」の形で、その組が共役な変数の組になる。
\end{rikai}
\end{mondai}

\mondaibangou{3}
\begin{mondai}{078-legendre}{計算}{Legendre 変換の計算，$x = -\dd g/\dd y$ と凸性の逆転，逆 Legendre 変換，平行移動の情報が保たれること}
\begin{teisuu}
  気体定数 & $R$ & \qty{8.314}{\joule\per\kelvin\per\mole}\\
\end{teisuu}
$x > 0$ で定義された関数 $f(x) = x\ln x - x$ を考える。
  \begin{shoumon}
    \item $f$ が下に凸であることを確かめ，Legendre 変換 $g(y)$ を \noteshiki{10.16} で求めよ。
    \item $g$ について \noteshiki{10.17} が成り立つこと，$g$ が上に凸であることを確かめよ。
    \item $g(y)$ を逆 Legendre 変換 \noteshiki{10.18} して $f(x)$ に戻ることを確かめよ。
    \item 定数 $C$ だけ平行移動した $f_{C}(x) = f(x - C)$（$x > C$）の Legendre 変換 $g_{C}(y)$ を求め，$g(y)$ との違いを述べよ。また，単純に $f_{C}(x)$ に $x = x_{C}(y)$ を代入しただけの関数と比べよ。
  \end{shoumon}

\begin{kotaeran}
  \machigai{4}{△}{計算} $x_{C}(y)=\ee^{y}+C$ と $g_{C}(y)=f_{C}(x_{C}(y))-yx_{C}(y)=y\ee^{y}-\ee^{y}-y(\ee^{y}+C)$ までは正しい。次の行で $-\ee^{y}+C$ としたが、$-y(\ee^{y}+C)=-y\ee^{y}-Cy$ なので $g_{C}(y)=-\ee^{y}-Cy=g(y)-Cy$ である。そのため「$g(y)$ を $-C$ 平行移動したもの」という結論も違う。単純に代入した $f_{C}(x_{C}(y))=(y-1)\ee^{y}$ が $C$ によらない、は正しい。
\end{kotaeran}
\begin{rikai}
  \item 平行移動の情報は、$g_{C}=g-Cy$ の傾き $-C$ の1次の項として残る。$f_{C}$ の接線（傾き $y$）は $f$ の接線を右に $C$ ずらしたもので、$x=0$ での切片は $-Cy$ だけ変わる。
  \item 検算：Legendre 変換の導関数の関係 $x=-\dd g/\dd y$、つまり $-g_{C}'(y)=x_{C}(y)$ が成り立つかを見る。$g_{C}=-\ee^{y}-Cy$ なら $-g_{C}'=\ee^{y}+C=x_{C}$ で合う。$-\ee^{y}+C$ では $-g_{C}'=\ee^{y}$ となって合わない。
\end{rikai}
\end{mondai}

\level{2}
\mondaibangou{4}
\begin{mondai}{079-power-law}{導出}{$U(S,V,n)$ の示量性・単調性・凸性からの制限（$C_{V}>0$，$\kappa_{S}>0$），Euler の関係式}
\begin{teisuu}
  気体定数 & $R$ & \qty{8.314}{\joule\per\kelvin\per\mole}\\
\end{teisuu}
一成分流体の内部エネルギーが
  \[ U(S,V,n) = KS^{a}V^{b}n^{c} \]
  の形で与えられるとする（$K > 0$，$S > 0$，$a$，$b$，$c$ は実数の定数）。
  \begin{shoumon}
    \item 内部エネルギーが示量変数であることから，$a$，$b$，$c$ が満たすべき条件を求めよ。
    \item $T > 0$ かつ $p > 0$ であることから，$a$，$b$ の符号を決めよ。$T$，$p$，$\mu$ を $U$，$S$，$V$，$n$ と $a$，$b$，$c$ で表せ。
    \item 定積比熱 $C_{V}$ と断熱圧縮率 $\kappa_{S}$ を \noteshiki{10.9}，\noteshiki{10.10} の右側の形で計算し，それらが正であるための $a$，$b$ の条件を求めよ。
    \item (1) の条件のもとで，Euler の関係式 \noteshiki{10.12} が成り立つことを示せ。Level 1 で扱った $U = KS^{3}/(nV)$ が，ここで求めた条件をすべて満たすことも確かめよ。
  \end{shoumon}

\begin{kotaeran}
  \machigai{3}{○}{記述} $(\partial T/\partial S)_{V,n}=a(a-1)U/S^{2}$、$(\partial p/\partial V)_{S,n}=-b(b-1)U/V^{2}$ と、条件 $a>1$、$b<0$ は正しい（$C_{V}=S^{2}T/(a(a-1)U)$ は $T=aU/S$ を入れると $S/(a-1)$ になる）。$C_{V}>0$、$\kappa_{S}>0$ が $U$ の凸性（$T$ が $S$ の非減少関数、$p$ が $V$ の非増加関数）から来る条件であることが書かれていない。
\end{kotaeran}
\begin{rikai}
  \item $C_{V}=T/(\partial T/\partial S)_{V,n}$ の分母は $\partial^{2}U/\partial S^{2}$、$\kappa_{S}$ の分母の $-(\partial p/\partial V)_{S,n}$ は $\partial^{2}U/\partial V^{2}$ である。$U$ が下に凸ならどちらも正で、$C_{V}>0$、$\kappa_{S}>0$ は「熱力学的な安定性」の条件である。
\end{rikai}
\end{mondai}

\mondaibangou{5}
\begin{mondai}{080-gibbs-duhem}{計算}{Gibbs-Duhem の式による化学ポテンシャルの圧力依存性，等温圧縮率を密度で表す式，第8章の状態方程式を使う}
\begin{teisuu}
  気体定数 & $R$ & \qty{8.314}{\joule\per\kelvin\per\mole}\\
\end{teisuu}
物質量 $n$ の一成分気体の状態方程式が，第8章で扱った
  \[ V = \frac{nRT}{p} + nb \]
  （$b$ は正の定数）で与えられるとする。
  \begin{shoumon}
    \item Gibbs-Duhem の式 \noteshiki{10.13} を使い，温度一定のもとで化学ポテンシャルの差 $\mu(T,p) - \mu(T,p_{0})$ を求めよ。
    \item 等温圧縮率 $\kappa_{T}$ を，定義から直接求めたものと，\noteshiki{10.15} の真ん中の形（密度 $\rho = n/V$ の圧力微分）で求めたものが一致することを確かめよ。
    \item $b = \qty{3.20e-5}{\metre\cubed\per\mole}$，$T = \qty{300}{\kelvin}$ のとき，圧力を \qty{1.00e5}{\pascal} から \qty{1.00e7}{\pascal} に上げたときの $\mu$ の変化を有効数字3桁で求めよ。そのうち $b$ による部分はいくらか。
  \end{shoumon}

\begin{kotaeran}
  \machigai{2}{○}{記述} 定義から求めた $\kappa_{T}$ と $(1/\rho)(\partial\rho/\partial p)_{T}$ が $RT/(p(bp+RT))$ で一致することは正しい。密度の微分で $n$ を固定する添字を外してよい理由（$\rho=n/V=p/(RT+bp)$ は示強変数で $T$ と $p$ だけで決まり、$n$ によらない）が書かれていない。
\end{kotaeran}
\begin{rikai}
  \item $V$ は示量変数なので $(\partial V/\partial p)_{T,n}$ と $n$ の固定が必要だが、$\rho$ や $v=V/n$ のような示強変数は $(T,p)$ だけの関数なので、$(\partial\rho/\partial p)_{T}$ と書ける。$\kappa_{T}=(1/\rho)(\partial\rho/\partial p)_{T}$ の形はこれを使っている。
\end{rikai}
\end{mondai}

\kaitouhen
\level{1}
\kaisetsubangou{2}
\begin{kaisetsu}{077-u-tv}
\houshin $U$ が $V$ によらないので，エネルギー方程式から $p$ の満たす偏微分方程式が得られる。それを解くと任意関数が残ることから，$U(T,V)$ が完全でないことが分かる。$U(S,V)$ では偏微分係数が直接 $T$ と $p$ を与えることと対比する。
\keisan \shoumonkai{1} $(\partial U/\partial V)_{T} = 0$ なので，\noteshiki{10.3} から $T(\partial p/\partial T)_{V,n} = p$。$V$，$n$ を固定すると
\[ \frac{\partial}{\partial T}\left(\frac{p}{T}\right) =\frac{1}{T}\left(\pdv{p}{T}\right)_{V,n} - \frac{p}{T^{2}} = 0 \]
なので，$p/T$ は $V$ と $n$ だけの関数 $f(V,n)$ であり，$p = Tf(V,n)$ となる。$p = nRT/V$ は $f = nR/V$，$p = nRT/V + bn^{2}T/V^{2}$ は $f = nR/V + bn^{2}/V^{2}$ の場合で，どちらも $T(\partial p/\partial T)_{V,n} = p$ を満たし，$U = ncT$ と矛盾しない。

\shoumonkai{2} $U(T,V)$ から状態方程式を求めるには，\noteshiki{10.3} を $p(T,V)$ についての微分方程式として解く必要があり，(1) のように任意関数（ここでは $Tf(V,n)$ の形の項）が残る。一方，$U(S,V)$ の偏微分は \noteshiki{10.2} で $T(S,V)$ と $p(S,V)$ そのものを与え，偏微分係数を含まない。2式から $S$ を消去すれば $p(T,V)$ が一意に決まる。熱力学関数が完全かどうかは変数の選び方に依存し，$U$ の自然な変数の組は $(S,V)$ である。

\shoumonkai{3} 物質の出入りがある準静的過程では $\delta W = -p\dd V + \mu\dd n$（\noteshiki{10.5}）であり，$\dd U = T\dd S - p\dd V + \mu\dd n$ に現れる共役な変数の組は $(T,S)$，$(p,V)$，$(\mu,n)$ である。化学ポテンシャルは，物質量を $\dd n$ だけ増やすときに（エントロピーと体積を一定に保って）系に与えられる仕事（化学的仕事）の係数である。
\kotae{\shoumonkai{1} $p = Tf(V,n)$。2つの状態方程式はどちらも $T(\partial p/\partial T)_{V,n} = p$ を満たす \shoumonkai{2} $U(S,V)$ の偏微分は $T$ と $p$ を直接与え，$S$ を消去して $p(T,V)$ が一意に決まる（$U$ の自然な変数は $(S,V)$） \shoumonkai{3} $\delta W = -p\dd V + \mu\dd n$。$(T,S)$，$(p,V)$，$(\mu,n)$。$\mu$ は化学的仕事 $\mu\dd n$ の係数}
\begin{ronten}
  \item \shoumonkai{1} $p/T$ の温度微分が0になることを示し，任意関数が残ることを述べる。
  \item \shoumonkai{2} $U(T,V)$ では偏微分係数を含む式しか得られないのに対し，$U(S,V)$ では偏微分が示強変数そのものになることを対比して書く。
  \item \shoumonkai{3} $\mu\dd n$ が物質の授受に伴う仕事であることを述べる。
\end{ronten}
\kenzan (1) の2つの状態方程式は $b \to 0$ で一致する。$p = nRT/V + bn^{2}T/V^{2}$ の追加の項も $T$ に比例するので，エネルギー方程式の右辺 $T(\partial p/\partial T) - p$ に寄与しない。
\tsumazuki 「$U$ が $T$ だけの関数なら理想気体」と決めつけない。状態方程式は $p = Tf(V,n)$ の形まで絞られるだけである。
\end{kaisetsu}

\kaisetsubangou{3}
\begin{kaisetsu}{078-legendre}
\houshin $y = f'(x)$ を $x$ について解き，\noteshiki{10.16} に代入する。逆変換は $x = -g'(y)$ を $y$ について解いて \noteshiki{10.18} に代入する。平行移動では $x(y)$ が $C$ だけずれることに注意する。
\keisan \shoumonkai{1} $f'(x) = \ln x + 1 - 1 = \ln x$，$f''(x) = 1/x > 0$ なので $f$ は下に凸である。$y = \ln x$ から $x(y) = \ee^{y}$。
\[ g(y) = f(\ee^{y}) - y\ee^{y} = y\ee^{y} - \ee^{y} - y\ee^{y} = -\ee^{y} \]
\shoumonkai{2} $-g'(y) = \ee^{y} = x(y)$ で \noteshiki{10.17} が成り立つ。$g''(y) = -\ee^{y} < 0$ なので $g$ は上に凸である。

\shoumonkai{3} $x = -g'(y) = \ee^{y}$ を $y$ について解くと $y(x) = \ln x$。
\[ g(y(x)) + xy(x) = -x + x\ln x = f(x) \]
\shoumonkai{4} $f_{C}'(x) = \ln(x - C)$ なので $x_{C}(y) = C + \ee^{y}$，$f_{C}(x_{C}(y)) = f(\ee^{y}) = y\ee^{y} - \ee^{y}$ である。
\[ g_{C}(y) = f_{C}(x_{C}(y)) - yx_{C}(y) = y\ee^{y} - \ee^{y} - y(C + \ee^{y}) = -\ee^{y} - Cy = g(y) - Cy \]
$g_{C}$ は $C$ を含み，平行移動の情報が保たれる。一方，単純に代入しただけの $f_{C}(x_{C}(y)) = y\ee^{y} - \ee^{y}$ は $C$ によらず，$C$ の情報が失われる。
\kotae{\shoumonkai{1} $f'' = 1/x > 0$，$g(y) = -\ee^{y}$ \shoumonkai{2} $-g'(y) = \ee^{y} = x$，$g'' = -\ee^{y} < 0$ \shoumonkai{3} $y(x) = \ln x$ で $g(y(x)) + xy(x) = x\ln x - x$ \shoumonkai{4} $g_{C}(y) = -\ee^{y} - Cy$（$C$ を含む）。単純な代入 $f_{C}(x_{C}(y)) = y\ee^{y} - \ee^{y}$ は $C$ によらない}
\kenzan (1) 幾何学的意味：$x = 1$ での接線は傾き $y = 0$，$f(1) = -1$ なので切片は $-1$ で，$g(0) = -\ee^{0} = -1$ と一致する。(4) で $C = 0$ なら $g_{C} = g$ に戻る。
\tsumazuki Legendre 変換の符号は要点整理の定義 $g = f - yx$ に従う（$yx - f$ と定義する流儀もあるので混ぜない）。
\end{kaisetsu}

\level{2}
\kaisetsubangou{4}
\begin{kaisetsu}{079-power-law}
\houshin 示量性は $\lambda$ 倍したときの指数の比較で，$T$，$p$，$\mu$ は \noteshiki{10.4} の偏微分で求める。$C_{V}$ と $\kappa_{S}$ は \noteshiki{10.9}，\noteshiki{10.10} の右側の形で計算し，正であるための条件を出す。
\keisan \shoumonkai{1} $U(\lambda S,\lambda V,\lambda n) = \lambda^{a+b+c}U(S,V,n)$ が任意の $\lambda > 0$ で $\lambda U$ に等しいので，$a + b + c = 1$。

\shoumonkai{2} べき関数なので
\[ T = \left(\pdv{U}{S}\right)_{V,n} = \frac{aU}{S}, \qquad p = -\left(\pdv{U}{V}\right)_{S,n} = -\frac{bU}{V}, \qquad \mu = \left(\pdv{U}{n}\right)_{S,V} = \frac{cU}{n} \]
$U > 0$，$S > 0$，$V > 0$ なので，$T > 0$ から $a > 0$，$p > 0$ から $b < 0$ である。

\shoumonkai{3} $(\partial T/\partial S)_{V,n} = a(a-1)KS^{a-2}V^{b}n^{c} = a(a-1)U/S^{2}$ なので
\[ C_{V} = T\Big/\left(\pdv{T}{S}\right)_{V,n} = \frac{aU}{S}\cdot\frac{S^{2}}{a(a-1)U} = \frac{S}{a-1} \]
$C_{V} > 0$ には $a > 1$ が必要である。また $(\partial p/\partial V)_{S,n} = -b(b-1)U/V^{2}$ なので
\[ \kappa_{S} = -1\Big/V\left(\pdv{p}{V}\right)_{S,n} = \frac{V}{b(b-1)U} \]
$\kappa_{S} > 0$ には $b(b-1) > 0$，すなわち $b < 0$ または $b > 1$ が必要で，(2) の $b < 0$ ならば満たされる。

\shoumonkai{4} (2) の結果から
\[ TS - pV + \mu n = aU + bU + cU = (a + b + c)U = U \]
で，(1) の条件のもとで Euler の関係式が成り立つ。Level 1 の $U = KS^{3}/(nV)$ は $a = 3$，$b = -1$，$c = -1$ の場合で，$a + b + c = 1$，$a = 3 > 1$，$b = -1 < 0$ をすべて満たす。このとき $C_{V} = S/2$ で，Level 1 で求めた $C_{V} = 3U/(2T) = 3S/6 = S/2$（$U = TS/3$）と一致する。
\kotae{\shoumonkai{1} $a + b + c = 1$ \shoumonkai{2} $a > 0$，$b < 0$。$T = aU/S$，$p = -bU/V$，$\mu = cU/n$ \shoumonkai{3} $C_{V} = S/(a-1)$，$\kappa_{S} = V/(b(b-1)U)$。$a > 1$（$b < 0$ なら $\kappa_{S} > 0$ は自動的） \shoumonkai{4} $TS - pV + \mu n = (a + b + c)U = U$。$a = 3$，$b = c = -1$ は条件を満たす}
\begin{ronten}
  \item \shoumonkai{1} 示量性を任意の $\lambda > 0$ について要求し，指数を比べることを書く。
  \item \shoumonkai{3} $C_{V} > 0$，$\kappa_{S} > 0$ が $U$ の凸性（$T$ が $S$ の非減少関数，$p$ が $V$ の非増加関数）から来る条件であることを述べる。
  \item \shoumonkai{4} Euler の関係式が示量性（1次の同次性）から出ることを述べる。
\end{ronten}
\kenzan $C_{V} = S/(a-1)$ はエントロピーの次元，$\kappa_{S} = V/(b(b-1)U)$ は体積／エネルギー＝$1/$圧力 の次元で正しい。$a \to 1$ では $T = U/S$ が $S$ によらなくなり，$C_{V} \to \infty$ になる。
\tsumazuki $C_{V}$ を $(\partial U/\partial T)_{V}$ で求めようとすると $U$ を $T$ で表し直す必要があり手間がかかる。\noteshiki{10.9} の $T/(\partial T/\partial S)_{V,n}$ の形を使う。
\end{kaisetsu}

\kaisetsubangou{5}
\begin{kaisetsu}{080-gibbs-duhem}
\houshin 温度一定では Gibbs-Duhem の式 \noteshiki{10.13} から $n\dd\mu = V\dd p$ となるので，状態方程式の $V/n$ を圧力で積分すれば $\mu$ の圧力依存性が求まる。等温圧縮率は定義と \noteshiki{10.15} の両方で計算して比べる。
\keisan \shoumonkai{1} $\dd T = 0$ では $\dd\mu = (V/n)\dd p$ で，$V/n = RT/p + b$ なので
\[ \mu(T,p) - \mu(T,p_{0}) = \int_{p_{0}}^{p}\left(\frac{RT}{p'} + b\right)\dd p' = RT\ln\frac{p}{p_{0}} + b(p - p_{0}) \]
\shoumonkai{2} 定義から，$(\partial V/\partial p)_{T,n} = -nRT/p^{2}$ なので
\[ \kappa_{T} = \frac{1}{V}\cdot\frac{nRT}{p^{2}} = \frac{nRT}{p^{2}}\cdot\frac{p}{n(RT + bp)} = \frac{RT}{p(RT + bp)} \]
一方，$\rho = n/V = p/(RT + bp)$ なので
\[ \left(\pdv{\rho}{p}\right)_{T} = \frac{(RT + bp) - bp}{(RT + bp)^{2}} = \frac{RT}{(RT + bp)^{2}}, \qquad \frac{1}{\rho}\left(\pdv{\rho}{p}\right)_{T} = \frac{RT}{p(RT + bp)} \]
で，両者は一致する。

\shoumonkai{3} $RT\ln(p/p_{0}) = 8.314 \times 300 \times \ln 100 = 2494.2 \times 4.6052 = \qty{11486}{\joule\per\mole}$，$b(p - p_{0}) = \num{3.20e-5} \times \num{9.90e6} = \qty{316.8}{\joule\per\mole}$ なので
\[ \Delta\mu = 11486 + 317 = \qty{1.18e4}{\joule\per\mole} \]
そのうち $b$ による部分は \qty{317}{\joule\per\mole} である。
\kotae{\shoumonkai{1} $\mu(T,p) - \mu(T,p_{0}) = RT\ln(p/p_{0}) + b(p - p_{0})$ \shoumonkai{2} どちらも $\kappa_{T} = RT/(p(RT + bp))$ \shoumonkai{3} $\Delta\mu = \qty{1.18e4}{\joule\per\mole}$，$b$ による部分 \qty{317}{\joule\per\mole}}
\begin{ronten}
  \item \shoumonkai{1} 温度一定で Gibbs-Duhem の式の $S\dd T$ が消えることを述べる。
  \item \shoumonkai{2} $\rho$ が示強変数で $n$ によらないので，$n$ 固定の条件を外して密度の圧力微分で書けることを述べる。
\end{ronten}
\kenzan $b \to 0$ で理想気体の $\mu - \mu_{0} = RT\ln(p/p_{0})$，$\kappa_{T} = 1/p$ に戻る。$\mu$ の次元は \unit{\joule\per\mole} で，$b\,p$（体積／物質量 × 圧力）も同じ次元である。$\kappa_{T} < 1/p$ で，体積 $nb$ の分だけ圧縮されにくい。
\tsumazuki $\mu$ は単位物質量当たりの量（示強変数）なので，$\dd\mu = V\dd p$ と $n$ を落とさない。
\end{kaisetsu}

\end{document}
